% !TEX program = uplatex
\documentclass[uplatex,dvipdfmx,a4paper,11pt]{jsarticle}

\usepackage[margin=22mm]{geometry}
\usepackage{amsmath,amssymb,amsthm,mathtools,bm,stmaryrd}
\usepackage{mathpartir}
\usepackage{booktabs,array,tabularx,longtable}
\usepackage{enumitem}
\def\pgfsysdriver{pgfsys-dvipdfmx.def}
\usepackage{graphicx}
\usepackage[most]{tcolorbox}
\usepackage{hyperref}
\usepackage{pxjahyper}
\usepackage{cleveref}
\hypersetup{
  colorlinks=true,
  linkcolor=blue!48!black,
  citecolor=blue!48!black,
  urlcolor=blue!48!black,
  pdftitle={Colored Weighted Simple-sub},
  pdfsubject={effect-only annotations, complete handlers, weighted constraints, and dynamic color elaboration},
  pdfauthor={}
}

\setlength{\parindent}{1em}
\setlength{\parskip}{0.22em}
\setlist[itemize]{topsep=0.35em,itemsep=0.15em,parsep=0pt}
\setlist[enumerate]{topsep=0.35em,itemsep=0.15em,parsep=0pt}
\allowdisplaybreaks

\definecolor{deepblue}{RGB}{38,74,118}
\definecolor{softblue}{RGB}{238,245,252}
\definecolor{softgreen}{RGB}{239,248,241}
\definecolor{softyellow}{RGB}{255,249,226}
\definecolor{softred}{RGB}{253,239,239}
\definecolor{softgray}{RGB}{246,247,249}

\newtcolorbox{plainbox}[1][]{
  breakable,colback=softblue,colframe=deepblue!70,boxrule=0.6pt,
  arc=1.2mm,left=2.5mm,right=2.5mm,top=1.7mm,bottom=1.7mm,#1}
\newtcolorbox{intuition}[1][]{
  breakable,colback=softgreen,colframe=green!45!black,boxrule=0.55pt,
  arc=1.2mm,title={直感},fonttitle=\bfseries,#1}
\newtcolorbox{readingbox}[1][]{
  breakable,colback=softyellow,colframe=orange!65!black,boxrule=0.55pt,
  arc=1.2mm,title={読み方},fonttitle=\bfseries,#1}
\newtcolorbox{statusbox}[1][]{
  breakable,colback=softred,colframe=red!55!black,boxrule=0.55pt,
  arc=1.2mm,title={証明状況},fonttitle=\bfseries,#1}
\newtcolorbox{formalbox}[1][]{
  breakable,colback=softgray,colframe=black!45,boxrule=0.5pt,
  arc=1.0mm,#1}

\newtheoremstyle{jpplain}
  {0.7em}{0.7em}{\normalfont}{}% above below body indent
  {\bfseries}{。}{0.5em}{}
\theoremstyle{jpplain}
\newtheorem{definition}{定義}[section]
\newtheorem{lemma}[definition]{補題}
\newtheorem{proposition}[definition]{命題}
\newtheorem{theorem}[definition]{定理}
\newtheorem{corollary}[definition]{系}
\newtheorem{conjecture}[definition]{予想}
\newtheorem{remark}[definition]{注意}
\newtheorem{example}[definition]{例}
\renewcommand{\proofname}{証明}

% common symbols
\newcommand{\Eff}{\mathsf{Eff}}
\newcommand{\Fam}{\mathsf{Fam}}
\newcommand{\Op}{\mathsf{Op}}
\newcommand{\Id}{\mathsf{Id}}
\newcommand{\Color}{\mathsf{Color}}
\newcommand{\Unit}{\mathbf 1}
\newcommand{\Nat}{\mathbb N}
\newcommand{\Int}{\mathbb Z}
\newcommand{\Pow}{\mathcal P}
\newcommand{\dom}{\mathsf{dom}}
\newcommand{\fv}{\mathsf{fv}}
\newcommand{\fresh}{\mathsf{fresh}}
\newcommand{\subty}{\mathrel{<:}}
\newcommand{\join}{\mathbin{\sqcup}}
\newcommand{\meet}{\mathbin{\sqcap}}
\newcommand{\eff}{\mathbin{!}}
\newcommand{\row}[2]{\{#1\,;\,#2\}}
\newcommand{\emptyrow}{\varnothing}
\newcommand{\starann}{\ast}
\newcommand{\gatom}[2]{\mathsf g_{#1}\!\{#2\}}
\newcommand{\patom}[1]{\mathsf p_{#1}}
\newcommand{\weight}{\mathsf W}
\newcommand{\zeroW}{\mathbf 0}
\newcommand{\coef}[2]{#1(#2)}
\newcommand{\capW}[2]{\mathsf{cap}_{#1}(#2)}
\newcommand{\exposeW}[2]{\mathsf{expose}_{#1}(#2)}
\newcommand{\wocc}[2]{#1^{\langle #2\rangle}}
\newcommand{\wsub}[4]{#1\ @#2\mathrel{<:}@#3\ #4}
\newcommand{\lb}{\mathsf{lb}}
\newcommand{\ub}{\mathsf{ub}}
\newcommand{\solve}{\mathsf{solve}}
\newcommand{\coalesce}{\mathsf{coal}}
\newcommand{\cleanup}{\mathsf{cleanup}}
\newcommand{\inst}{\mathsf{inst}}
\newcommand{\gen}{\mathsf{gen}}
\newcommand{\Pol}{\mathsf{Pol}}
\newcommand{\CovG}{\mathsf{CovGuard}}
\newcommand{\CovP}{\mathsf{CovPaint}}
\newcommand{\Live}{\mathsf{Live}}
\newcommand{\Elab}{\mathcal E}
\newcommand{\erase}{\mathsf{erase}}
\newcommand{\flow}{\Phi}
\newcommand{\colors}{\kappa}
\newcommand{\visible}{\mathsf{visible}}
\newcommand{\Guard}{\mathcal G}
\newcommand{\Paint}{\mathcal P}
\newcommand{\ret}{\mathsf{return}}
\newcommand{\oper}{\mathsf{op}}
\newcommand{\thunk}{\mathsf{thunk}}
\newcommand{\force}{\mathsf{force}}
\newcommand{\perform}{\mathsf{perform}}
\newcommand{\scope}{\mathsf{scope}}
\newcommand{\getid}{\mathsf{get\_id}}
\newcommand{\pushid}{\mathsf{push\_id}}
\newcommand{\popid}{\mathsf{pop\_id}}
\newcommand{\guard}{\mathsf{guard}}
\newcommand{\paint}{\mathsf{paint}}
\newcommand{\poppaint}{\mathsf{pop\_paint}}
\newcommand{\handle}{\mathsf{handle}}
\newcommand{\with}{\mathsf{with}}
\newcommand{\raw}{\mathsf{raw}}
\newcommand{\markg}{\mathsf{markGuard}}
\newcommand{\markp}{\mathsf{markPaint}}
\newcommand{\seqeff}{\mathsf{seq}}
\newcommand{\annok}{\mathsf{ann}}
\newcommand{\later}{\triangleright}

\crefname{definition}{定義}{定義}
\crefname{lemma}{補題}{補題}
\crefname{proposition}{命題}{命題}
\crefname{theorem}{定理}{定理}
\crefname{corollary}{系}{系}
\crefname{conjecture}{予想}{予想}
\crefname{example}{例}{例}

\title{Colored Weighted Simple-sub\\
\large effect-only skeleton・明示的 ref・型指向 force の形式化（議論稿）}
\author{}
\date{2026年6月21日}

\begin{document}
\pagestyle{plain}
\maketitle

\begin{abstract}
本稿は、値型を通常どおり推論しつつ、ユーザが computation slot の effect 境界だけを
注釈する小言語を定める。注釈は effect-only skeleton
\[
  K ::= q \mid K\to K,
  \qquad
  q ::= \emptyrow \mid \starann \mid H
\]
であり、値型を直接指定しない。表面構文上の省略は \(\starann\) の糖衣であって、
推論規則は省略形を扱わない。

項の側では、束縛名を値として使う操作を \(\mathsf{ref}\ x\) と明示する。
また、effectful computation は thunk として値に閉じ込められ、実行は \(\force\) によって行う。
表面構文では \(\force\) を省けるが、core calculus では型に従って挿入された \(\force\) を明示する。
この分離により、静的 id を発行する場所（\(\mathsf{ref}\)）と effect が実際に流れる場所
（\(\force\)）を区別できる。

型推論の中心は Simple-sub 型の上下界 solver であり、型比較を
\(\wsub{T}{L}{R}{U}\) という左右に重みを持つ不等式へ拡張する。
重みは静的識別子ごとの符号付き加法で、正の生成子を \emph{guard}、負の生成子を
\emph{paint} と呼ぶ。handler が effect row を分解するとき、重みを直接 row の tail へ書き込まず、
fresh な中間 tail を一度挟んでから重み付き不等式を replay する。

静的な guard/paint と、実行時の push/pop は区別する。型解決後の elaboration は
\(\getid\)、\(\guard(c,H,v)\)、\(\pushid(c)\)、\(\popid(c)\)、\(\paint(c,v)\)
を挿入する。同一の静的 id から生じる guard と push には同じ動的 color id を渡す。
関数を起動すると id を動的 stack に積み、通常 return でも effect による脱出でも、関数を抜けると
id を外して、境界を越える値・operation payload・継続に paint を付ける。
したがって color は値になるまで粘着的に伝播する。

各 effect family は一つの operation だけを持ち、handler はその family を完全に処理するものだけを許す。
let 多相では型変数、row 変数に加えて静的 id も量化し、instantiate ごとに fresh な id と
\(\getid\) を生成する。再帰は本体を調べる前に effect-only skeleton から自己型の下界を置く。
一般化時の cleanup は、共変 guard と共変 paint の双方を持つ id だけを観測可能として残す。
本文では局所的な代数則、weighted row split、動的 color hygiene を証明し、solver 全体の完全性・主型性、
cleanup の文脈同値性、elaboration の型保存を明示的な予想として分離する。
\end{abstract}

{\small\tableofcontents}

\section{一ページで見る全体像}

\begin{plainbox}[title={設計の八段階},fonttitle=\bfseries]
\begin{enumerate}[label=\textbf{\arabic*.}]
  \item \textbf{注釈は effect-only skeleton である。}
  ユーザは値型を書かず、各 computation slot に
  \[
    q\in\{\emptyrow,\starann,H\}
  \]
  だけを書く。表面構文上の省略は \(\starann\) の糖衣であり、推論規則には現れない。

  \item \textbf{表示型は普通の型である。}
  例えば skeleton から作られる引数 $f$ の環境内の型は
  \[
    f : \alpha\eff\emptyrow \to \beta\eff S
  \]
  であり、$S$ 自体は重み付き型ではない。

  \item \textbf{\(\mathsf{ref}\ x\) が値参照を表す。}
  束縛名 $x$ を値として取り出すには必ず \(\mathsf{ref}\ x\) と書く。
  $f$ の参照出現だけが
  \[
    \alpha\eff\emptyrow \to \beta\eff\wocc{S}{\gatom a H}
  \]
  を得る。型変数そのものと、その使用経路を区別する。

  \item \textbf{\(\force\) が effect thunk を実行する。}
  表面構文では \(\force\) を省けるが、core では型に従って挿入された \(\force\) を見る。
  したがって id 発行の候補地点は \(\mathsf{ref}\)、effect が実際に流れる地点は \(\force\) である。

  \item \textbf{比較の左右に重みが乗り、bounds replay で加算される。}
  関数の反変位置では左右が交換される。変数の lower/upper bound をつなぐと、経路上の重みが足される。

  \item \textbf{handler は row を直接書き換えない。}
  \[
    \wsub{S}{\gatom a H}{\zeroW}{\row h T}
  \]
  から直接 $S=\row h{\wocc T{\gatom a H}}$ とはしない。fresh な $U$ を挟んで
  \[
    S\subty\row h U,
    \qquad
    \wsub{U}{\gatom a H}{\zeroW}{T}
  \]
  を生成する。

  \item \textbf{一般化時にだけ cleanup する。}
  共変 guard がなければ対応 paint は観測不能であり、共変 paint がなければ guard も外へ暴れられない。
  両方を持つ id だけを scheme に残す。

  \item \textbf{解決した型から core term を elaboration する。}
  retained id ごとに $c=\getid()$ を生成し、同じ $c$ を guard と push に渡す。
  関数脱出時には pop と paint を一緒に行い、値・payload・continuation に color を残す。
\end{enumerate}
\end{plainbox}

\begin{intuition}
静的重みは runtime stack そのものではない。動的 stack の順序と well-bracketedness は push/pop が担う。
静的重みは、その動的操作が型変数の流れに残す「符号付き残量」だけを可換な加法として記録する。
この分離により、型推論は Simple-sub の bounds graph と極性解析を使え、意味論は本物の id stack を使える。
\end{intuition}

\paragraph{中心例：compose}
次の式を考える。
\begin{verbatim}
let compose =
  lambda (f : {} -> *).
  lambda (g : *  -> *).
  lambda (x : {}).
    ref g (ref f (ref x))
\end{verbatim}
表面構文では \(\force\) を省けるが、core では型に従って概念的に
\[
  \force\bigl((\mathsf{ref}\ g)\,
    (\force((\mathsf{ref}\ f)(\mathsf{ref}\ x)))\bigr)
\]
へ elaboration される。
$f$ 自体の型は普通の $\alpha\eff\emptyrow\to\beta\eff S$ である。
しかし $f\,x$ の使用出現は fresh id $a$ の guard を持ち、$g$ の引数型へ流れる。
全体の共変出力には対応 paint が現れる。したがって cleanup 後にも
\[
 (\alpha\eff\emptyrow\to\beta\eff S)
 \to
 (\beta\eff\wocc S{\gatom a\Eff}\to\gamma\eff T)
 \to
 \alpha\eff\emptyrow\to\gamma\eff\wocc T{\patom a}
\]
が残る。$g$ の引数は $g$ に対して反変、$g$ 自体は compose に対して反変なので、
その guard は全体では共変位置にある。ここに Simple-sub の極性解析が直接効く。

\begin{statusbox}
この文書は「次に議論するための形式化」である。型・重み・handler・一般化・elaboration・動的 color の
接続を一つの体系として書く一方、完全な機械検証済み証明を装わない。
局所補題は証明し、全 solver の停止性・完全性・主型性、および cleanup 後の core term 消去は
証明計画を伴う予想として明示する。
\end{statusbox}

\section{表面言語と effect-only 注釈}

\subsection{Effect signature}

\begin{definition}[Effect signature]
$\Fam$ を effect family の集合とする。本稿では各 family $h\in\Fam$ に operation を一つだけ割り当てる。
\[
  \Sigma(h)=P_h\rightsquigarrow Q_h.
\]
$P_h$ は operation payload の値型、$Q_h$ は continuation を再開するときに渡す値型である。
同じ family に複数 operation をまとめる体系は、各 operation を別 family と見なすことで近似できる。
このため core calculus では incomplete handler を導入しない。
\end{definition}

\subsection{Effect-only skeleton}

\begin{definition}[effect slot 注釈]
core calculus の effect slot 注釈を
\[
 q ::= \emptyrow \mid \starann \mid H
 \qquad(H\subseteq_{\mathrm{fin}}\Fam)
\]
で定める。$\emptyrow$ は concrete empty boundary、$\starann$ は無制限、$H$ は具体的 family 集合である。
表面では $\emptyrow$ を \texttt{\{\}}、有限集合 $\{h_1,\ldots,h_n\}$ を
\texttt{\{h1, ..., hn\}} と書く。
表面構文上は注釈を省略してよいが、それは常に $\starann$ の糖衣である。
したがって推論規則は「省略された注釈」という場合分けを持たない。
\end{definition}

値型は書かず、effect slot だけを記述する skeleton を定める。
\[
\begin{array}{rcl}
 K &::=& q \mid K\to K.
\end{array}
\]
これは型そのものではなく、computation slot に置く effect boundary の形だけを指定する kind 的な構造である。
したがって
\[
  \emptyrow\to\{\mathsf{io},\mathsf{nondet}\}
\]
は値型を指定せず、domain と codomain の二つの computation slot だけを指定する。
推論時には fresh な値型と row 変数を立て、概念的には
\[
  A\eff\emptyrow \to B\eff\row{\{\mathsf{io},\mathsf{nondet}\}}S
\]
のような ordinary 型と、後述する occurrence transformer を生成する。

\begin{readingbox}
「共変注釈は高階関数を無視する」とは、ある computation slot $A\eff R$ に付いた注釈が
$R$ 上の重みだけを検査し、値型 $A$ の内部へ勝手に再帰しない、という意味である。
ただし skeleton に高階関数内部の注釈を明示した場合、その明示した各 slot は独立に翻訳される。
\end{readingbox}

\subsection{項構文}

議論の対象を次に限定する。
\[
\begin{array}{rcl}
 v &::=& () \mid \lambda(x:K):q.\ e \mid \thunk(e),\\
 e &::=& \ret\,v \mid \mathsf{ref}\ x \mid e\ e \mid \force(e) \mid \perform_h(e)\\
   && \mid\ \mathsf{handle}\;e\;\mathsf{with}\;h(y,k)\Rightarrow e_h\mid \ret\,x\Rightarrow e_r\\
   && \mid\ \mathsf{let}\ x=e_1\ \mathsf{in}\ e_2\\
   && \mid\ \mathsf{let\ rec}\ f:K=e_1\ \mathsf{in}\ e_2.
\end{array}
\]
束縛名 $x$ は、式位置でそのまま値を表さない。値を現在の computation へ取り出す操作は必ず
\(\mathsf{ref}\ x\) と書く。したがって静的 id と occurrence transformer が発火しうる場所も
\(\mathsf{ref}\) に限られる。
\(\lambda(x:K).\ e\) という表面構文は、返り slot を省いた
\(\lambda(x:K):\starann.\ e\) の糖衣である。
\(\force\) は core calculus の構文である。表面言語では、期待される型が computation であり、
手元の値が effect thunk であると分かる場所に \(\force\) を自動挿入してよい。
本稿の制約生成は、この型指向挿入後の core term を対象にする。
handler は family $h$ の唯一の operation を完全に処理し、operation branch と return branch を一つずつ持つ。
通常の評価順序に伴う effect の合流は $\seqeff$ で表し、本稿の color 機構と直交する部分は簡略化する。

\section{型・row・静的重み}

\subsection{表示型と内部型}

\begin{definition}[型]
値型、computation 型、effect row を次で定める。
\[
\begin{array}{rcl}
 A &::=& \alpha \mid \Unit \mid C\to C \mid \mu\alpha.A,\\
 C &::=& A\eff R,\\
 R &::=& \rho \mid \emptyrow \mid \row H R.
\end{array}
\]
$\row H R$ は有限 head $H$ と tail $R$ を持つ row pattern である。
本稿の小体系では head と tail は disjoint であるとし、$H$ の family は tail に再出現しない。
内部推論では型・row を cyclic graph として持ち、表示時に必要なら $\mu$ を導入する。
\end{definition}

\begin{remark}[表示型と thunk 表現]
上の \(C_1\to C_2\) は、関数が argument computation slot \(C_1\) から result computation slot \(C_2\) を作る、
という表示型である。core の実行表現では、effectful result は \(\thunk(e)\) として値に閉じ込められ、
\(\force\) によって実行される。この operational thunk は、ユーザが書く effect-only skeleton の一部ではなく、
型指向 elaboration の産物として扱う。
\end{remark}

重みは表示型そのものではなく、型または row の\emph{出現}に付く。
\[
  \wocc T w
\]
と書く。$T$ という型変数の全ての出現が重くなるわけではない。

\subsection{静的 id と signed weight}

\begin{definition}[静的 id 宣言]
静的 id context $\Delta$ は有限写像
\[
  \Delta : \Id \rightharpoonup \Pow(\Fam)
\]
である。$\Delta(a)=H$ は id $a$ の guard が開放する family 集合を表す。
同じ $a$ に異なる $H$ を割り当てない。
\end{definition}

\begin{definition}[重み]
重み $w$ は有限台を持つ整数値写像
\[
  w:\Id\to_{\mathrm{fin}}\Int
\]
である。加算は点ごとの整数加算、単位元は零写像 $\zeroW$ とする。
\[
  (w_1+w_2)(a)=w_1(a)+w_2(a).
\]
$w(a)=+1$ の生成子を
\[
  \gatom a{\Delta(a)}=\#a\{\Delta(a)\}(1)
\]
と書き \emph{guard} と呼ぶ。$w(a)=-1$ の生成子を
\[
  \patom a=\#a(\_)(-1)
\]
と書き \emph{paint} と呼ぶ。
\end{definition}

代数的には $(\weight,+,\zeroW)$ は可換群である。ただし本稿では、負の生成子を自由な逆操作と解釈せず、
well-bracketed な動的 push/pop から射影された signed monoid evidence と読む。

\begin{lemma}[重みの正規形]
任意の有限な guard/paint 列は、各 $a$ の整数係数 $w(a)$ を集めた一意な正規形を持つ。
\end{lemma}
\begin{proof}
生成子を id ごとに分類し、guard の個数から paint の個数を引く。整数加算の結合律・可換律と
加法逆元の一意性から、分類順序によらず同じ有限台写像を得る。
\end{proof}

\begin{definition}[handler capability と注釈露出]
正の balance を持つ id だけが handler の subtraction 許可を与える。
\[
 \capW\Delta w
 = \bigcap_{a\,:\,w(a)>0}\Delta(a),
\]
空の共通部分は $\Fam$ とする。
一方、共変注釈が検査する active guard 集合を
\[
 \exposeW\Delta w
 = \bigcup_{a\,:\,w(a)>0}\Delta(a)
\]
とする。具体注釈 $H$ は
\[
  \annok_\Delta(w,H) \quad\Longleftrightarrow\quad
  \exposeW\Delta w\subseteq H
\]
を要求する。ordinary row 自体には $R\subty H$ を課さない。
\end{definition}

\begin{example}
$w=\gatom a{\{\mathsf{write}\}}+\patom b$ なら、$b$ の paint は capability を与えず、
\[
 \capW\Delta w=\{\mathsf{write}\},\qquad
 \exposeW\Delta w=\{\mathsf{write}\}.
\]
返り注釈 $\{\mathsf{io}\}$ が要求するのは
$\{\mathsf{write}\}\subseteq\{\mathsf{io}\}$ であり、row tail $S\subseteq\{\mathsf{io}\}$ ではない。
\end{example}

\section{Weighted Simple-sub}

\subsection{重み付き不等式}

型比較を
\[
  \wsub{T}{L}{R}{U}
\]
と書く。$L$ は左の型出現までに通った重み、$R$ は右の型出現までに通った重みである。
型内の weighted occurrence に到達したら、重みを比較側へ移す。
\[
\inferrule*[right=W-Left]
 {\wsub{T}{w+L}{R}{U}}
 {\wsub{\wocc T w}{L}{R}{U}}
\qquad
\inferrule*[right=W-Right]
 {\wsub{T}{L}{w+R}{U}}
 {\wsub{T}{L}{R}{\wocc U w}}.
\]

構造型は通常どおり分解する。
\[
\inferrule*[right=W-Fun]
 {\wsub{D_1}{R}{L}{C_1}\\
  \wsub{C_2}{L}{R}{D_2}}
 {\wsub{C_1\to C_2}{L}{R}{D_1\to D_2}}
\]
\[
\inferrule*[right=W-Comp]
 {\wsub{A}{L}{R}{B}\\
  \wsub{S}{L}{R}{T}}
 {\wsub{A\eff S}{L}{R}{B\eff T}}.
\]
関数 domain で左右が交換されるため、重みの極性も自動的に反転する。

\subsection{重み付き bounds graph}

型変数状態 $\sigma$ は各変数 $\alpha$ に weighted lower/upper bounds を保存する。
\[
\begin{aligned}
 \lb_\sigma(\alpha)&\subseteq \mathsf{Type}\times\weight\times\weight,\\
 \ub_\sigma(\alpha)&\subseteq \mathsf{Type}\times\weight\times\weight.
\end{aligned}
\]
$\wsub{S}{L}{R}{\alpha}$ を lower bound、$\wsub{\alpha}{L}{R}{T}$ を upper bound として記録する。
新しい bound を追加したら、反対側の既存 bound と replay する。
\[
\inferrule*[right=W-Replay]
 {(S,L_s,R_s)\in\lb_\sigma(\alpha)\\
  (T,L_t,R_t)\in\ub_\sigma(\alpha)}
 {\wsub{S}{L_s+L_t}{R_s+R_t}{T}}
\]
この加算が、比較経路に沿って guard/paint が伝播する場所である。

\begin{formalbox}[title={変数ケースの実装的不変条件},fonttitle=\bfseries]
新しい bound は replay より先に状態へ挿入する。再帰的 replay が同じ変数へ戻ったとき、新しい bound が見える必要がある。
また比較 cache は型対だけでなく左右の正規化済み重みも含める。これらは Simple-sub の mutable bounds graph と同じ注意点である。
\end{formalbox}

\subsection{完全 weighted row split}

本体系の中心規則を与える。handler などから
\[
  \wsub{\alpha}{L}{R}{\row H\beta}
\]
が要求されたとする。$w=L+R$ と置く。

\begin{definition}[weighted row split]
\begin{enumerate}
  \item $w=\zeroW$ なら、通常の upper bound
  $\alpha\subty\row H\beta$ を登録する。
  \item $w\neq\zeroW$ なら、完全 handler 条件
  \[
    H\subseteq\capW\Delta w
  \]
  を要求する。満たさなければ型エラーである。
  \item fresh な row 変数 $\gamma$ を取り、
  \[
    \alpha\subty\row H\gamma,
    \qquad
    \wsub{\gamma}{w}{\zeroW}{\beta}
  \]
  を登録する。
\end{enumerate}
\end{definition}

\begin{readingbox}
直接 $\alpha\subty\row H{\wocc\beta w}$ としない。
一旦普通の row upper bound $\row H\gamma$ を作り、fresh tail $\gamma$ から期待 tail $\beta$ へ
重み付き edge を張る。この中間変数が bounds replay と共起簡約の接続点になる。
\end{readingbox}

\begin{lemma}[row split の局所健全性]
weighted row split が成功し、生成された二制約が成り立つなら、元の要求
$\wsub{\alpha}{L}{R}{\row H\beta}$ を満たす。
\end{lemma}
\begin{proof}
$w=0$ は定義そのもの。$w\neq0$ では
$\alpha\subty\row H\gamma$ により、$\alpha$ の値は head $H$ と tail $\gamma$ に分解される。
$\wsub{\gamma}{w}{0}{\beta}$ は元の比較経路の重みを tail に保存する。
row constructor の tail 共変性と推移律から、元の weighted comparison を再構成できる。
完全性条件 $H\subseteq\capW\Delta w$ は、この分解で handler が guard の許していない family を消費しないことを保証する。
\end{proof}

\begin{remark}
以前の一般仕様では、部分的にしか処理できない family を $H\cap\capW\Delta w$ として残す近似規則も考えられる。
本稿では各 family に operation を一つだけ割り当て、完全 handler だけを許すため、その分岐を削除する。
複数 operation を一 family にまとめる方式は、この完全体系に対する近似と見なせる。
\end{remark}

\subsection{極性付き coalescing}

推論後、Simple-sub と同様に正の変数出現は lower bounds、負の変数出現は upper bounds から展開する。
関数 domain で極性を反転し、weighted edge の重みを展開先の変数出現へ加算する。
概略は次である。
\[
\begin{aligned}
 \coalesce_\sigma^+(\alpha)
   &= \alpha\join
      \bigsqcup_{(T,L,R)\in\lb_\sigma(\alpha)}
        \coalesce_\sigma^+(\wocc T{L+R}),\\
 \coalesce_\sigma^-(\alpha)
   &= \alpha\meet
      \bigsqcap_{(T,L,R)\in\ub_\sigma(\alpha)}
        \coalesce_\sigma^-(\wocc T{L+R}).
\end{aligned}
\]
循環は polar な $\mu$ 型として閉じる。実装では compact graph を作ってから共起解析を行う。

\begin{intuition}
$f$ の型変数 $S$ に guard を直接保存するのではない。$f$ の参照出現から生じた edge に guard が乗り、
その edge が application、function variance、handler split を通って bounds graph を移動する。
最後に coalescing したときだけ、どの極性に重みが露出したかが表示型に現れる。
\end{intuition}

\section{注釈翻訳と項の制約生成}

\subsection{注釈の二つの役割}

concrete effect 注釈 $H$ は極性で役割が変わる。

\paragraph{反変 slot}
raw row 変数 $\rho$ と fresh id $a$ を作り、$\Delta(a)=H$ とする。
環境に保存する raw 型は $A\eff\rho$ のままである。一方、その値を \(\mathsf{ref}\) で参照した出現には
$\gatom aH$ を付ける occurrence transformer を保存する。
$H=\emptyrow$ なら、許可 family が空の guard として振る舞う。
$q=\starann$ は id を発行せず、transformer も恒等写像にする。

\paragraph{共変 slot}
raw row には制約を置かず、その slot に到達した重み $w$ だけについて
$\annok_\Delta(w,H)$ を検査する。値型内部へは再帰しない。
$q=\starann$ は検査を行わない。

\begin{formalbox}[title={環境 entry},fonttitle=\bfseries]
環境 entry を
\[
  \Gamma(x)=\forall\bar\alpha\,\bar\rho\,\bar a.\ (C;\Xi)
\]
と書く。$C$ は ordinary type scheme、$\Xi$ は effect slot ごとの occurrence transformer である。
lookup は scheme と id を freshen した後、$\Xi$ をその\emph{参照出現}に適用する。
したがって環境中の $f$ は普通の型を持ち、$\mathsf{ref}\ f$ だけが重くなる。
\end{formalbox}

\subsection{代表的な制約生成規則}

判断を
\[
  \Delta;\Gamma\vdash e\Rightarrow C\mid\mathcal C;M^\circ
\]
と書く。$C$ は raw inferred type、$\mathcal C$ は weighted constraints、$M^\circ$ は静的 id marker を持つ中間項である。

\paragraph{変数参照}
$\Gamma(x)=\sigma$ を instantiate し、ordinary type $C$ と transformer $\Xi$ を得る。
\[
\inferrule*[right=I-Ref]
 {\inst(\sigma)=(C,\Xi)}
 {\Delta;\Gamma\vdash \mathsf{ref}\ x\Rightarrow \Xi(C)\mid\emptyset;\mathsf{ref}\,x}
\]
裸の $x$ にはこの規則を適用しない。これにより、id 発行と occurrence transformer の適用地点は構文上明確になる。

\paragraph{lambda}
引数 skeleton を外側の関数 domain、すなわち負極性で翻訳し、raw 型 $C_x$、transformer $\Xi_x$、
fresh id 集合 $A_x$ を得る。本体を推論した raw result を $B\eff R$、その outer row の重みを $w_R$ とする。
具体返り注釈 $H_o$ があれば $\annok_\Delta(w_R,H_o)$ を検査する。その後、$A_x$ の各 id に対応する paint を
exported result occurrence に加える。
\[
\inferrule*[right=I-Lam]
 {\mathsf{neg}(K)=(C_x,\Xi_x,A_x)\\
  \Delta;\Gamma,x:(C_x;\Xi_x)\vdash e\Rightarrow B\eff\wocc R{w_R}\mid\mathcal C;M^\circ\\
  \mathsf{checkPos}(q,w_R)\\
  p_{A_x}=\sum_{a\in A_x}\patom a}
 {\Delta;\Gamma\vdash \lambda(x:K):q.e
  \Rightarrow ((C_x\to B\eff\wocc R{w_R+p_{A_x}})\eff\emptyrow)
  \mid\mathcal C;\lambda x.M^\circ}
\]
重要なのは、返り注釈の検査が paint を加える前に行われることである。

\paragraph{application}
$e_1$、$e_2$ を推論し、fresh な $B,R_i,R_o$ を取る。関数型比較を solver へ渡す。
通常の評価 effect は $\seqeff$ で合流する。表面構文で \(\force\) が省かれていた場合、
この規則へ入る前に、期待型に従って必要な \(\force\) が挿入されているものとする。
\[
\inferrule*[right=I-App]
 {\Gamma\vdash e_1\Rightarrow A_1\eff E_1\\
  \Gamma\vdash e_2\Rightarrow A_2\eff E_2\\
  \wsub{A_1}{0}{0}{(A_2\eff R_i)\to(B\eff R_o)}}
 {\Gamma\vdash e_1\ e_2\Rightarrow B\eff\seqeff(E_1,E_2,R_o)}
\]
重みは $A_1$、$A_2$、$R_o$ の出現に既に付いており、構造分解と replay で移動する。

\paragraph{operation}
\[
\inferrule*[right=I-Perform]
 {\Gamma\vdash e\Rightarrow P_h\eff E}
 {\Gamma\vdash \perform_h(e)\Rightarrow Q_h\eff\seqeff(E,\row{\{h\}}\emptyrow)}
\]
引数型の細部は本稿の color 機構と独立なので省略した。

\section{完全 handler}

\subsection{型付け規則}

scrutinee を
\[
  \Gamma\vdash e\Rightarrow B\eff\wocc S w
\]
と推論したとする。family $h$ を処理するため fresh tail $T$ を取り、solver に
\[
  \wsub{S}{w}{\zeroW}{\row{\{h\}}T}
\]
を要求する。weighted row split はさらに fresh $U$ を作り、
\[
  S\subty\row{\{h\}}U,
  \qquad
  \wsub{U}{w}{\zeroW}{T}
\]
を登録する。

operation branch では
\[
  y:P_h,
  \qquad
  k:Q_h\eff\emptyrow\to B\eff S
\]
とする。continuation の返り row は residual $T$ ではなく、元の scrutinee row $S$ である。
これは shallow continuation である。

branch と return branch を
\[
 \Gamma,y:P_h,k:Q_h\eff\emptyrow\to B\eff S\vdash e_h\Rightarrow C\eff E_h,
\]
\[
 \Gamma,x:B\vdash e_r\Rightarrow C\eff E_r
\]
と推論したとき、handler 全体は
\[
 C\eff\seqeff(T,E_h,E_r)
\]
を持つ。$T$ に到達する重みは、上の $U@w<:T$ と bounds replay によって自動的に伝わる。

\begin{formalbox}[title={Handler rule（概略）},fonttitle=\bfseries]
\[
\inferrule*[right=I-Handle]
 {\Gamma\vdash e\Rightarrow B\eff\wocc S w\\
  \wsub{S}{w}{0}{\row{\{h\}}T}\\
  \Gamma,y:P_h,k:Q_h\eff\emptyrow\to B\eff S\vdash e_h\Rightarrow C\eff E_h\\
  \Gamma,x:B\vdash e_r\Rightarrow C\eff E_r}
 {\Gamma\vdash \mathsf{handle}\;e\;\mathsf{with}\;h(y,k)\Rightarrow e_h\mid\mathsf{return}\;x\Rightarrow e_r
  \Rightarrow C\eff\seqeff(T,E_h,E_r)}
\]
\end{formalbox}

\subsection{小さい手計算：中間 tail を省かない}

次の wrapper を考える。branch と return branch は純粋に unit を返すとする。
\begin{verbatim}
lambda (f : {} -> {write}) : *.
  handle ref f () with
  | write(y, k) -> ()
  | return x    -> ()
\end{verbatim}
core では scrutinee の実行点に \(\force\) が挿入され、概念的には
\[
  \mathsf{handle}\;
    \force((\mathsf{ref}\ f)\,())
  \;\mathsf{with}\;\cdots
\]
を型付けする。

\begin{enumerate}
  \item 環境中の raw 型は
  \[
    f:\Unit\eff\emptyrow\to B\eff S.
  \]
  \item $f$ の参照出現だけが
  \[
    \Unit\eff\emptyrow\to B\eff\wocc S{\gatom a{\{\mathsf{write}\}}}
  \]
  となる。
  \item handler 比較は
  \[
    \wsub{S}{\gatom a{\{\mathsf{write}\}}}{0}{\row{\{\mathsf{write}\}}T}
  \]
  である。
  \item solver は直接 $S\subty\row{\{\mathsf{write}\}}{\wocc T{\gatom aH}}$ とせず、fresh $U$ を置いて
  \[
    S\subty\row{\{\mathsf{write}\}}U,
    \qquad
    \wsub{U}{\gatom a{\{\mathsf{write}\}}}{0}{T}
  \]
  とする。
  \item continuation は
  \[
    k:\Unit\eff\emptyrow\to B\eff S
  \]
  である。
  \item lambda の exported output に $\patom a$ が加わる。coalescing すると handler result の tail 経路には
  $\gatom aH+\patom a=0$ が現れる。
  \item 共変 paint が正規化で消えた後、その id に live な共変 paint が残らないため、input 側の guard も cleanup で消える。
\end{enumerate}
最終表示型は
\[
  (\Unit\eff\emptyrow\to B\eff\row{\{\mathsf{write}\}}T)
  \to \Unit\eff T
\]
となる。普通の handler 型に見えるが、その導出途中では weighted edge と fresh $U$ が必要である。

\section{let 多相・再帰・cleanup}

\subsection{Scheme と instantiate}

scheme は ordinary type variables、row variables、static ids を全て量化する。
\[
  \sigma ::= \forall\bar\alpha\,\bar\rho\,\bar a.\ \mathcal C\Rightarrow(C;\Xi;\Delta|_{\bar a}).
\]
instantiate は同じ scheme 内の同じ id を同じ fresh id へ写し、別 instantiate とは共有しない。
\[
  \inst(\bar a)=\bar a'\qquad(\bar a'\ \text{fresh}).
\]
これが runtime の fresh $\getid$ と対応する。

\subsection{let 多相}

$e_1$ を一段高い level で推論し、外側環境に逃げていない変数と id を一般化する。
一般化前に coalescing と cleanup を行う。値制限を採る場合は syntactic value だけを一般化するが、
本稿の color 機構はその選択と独立である。

\subsection{再帰}

\[
 \mathsf{let\ rec}\ f:K=e_1\ \mathsf{in}\ e_2
\]
では、まず skeleton $K$ を翻訳して raw skeleton $C_s$ と transformer $\Xi_s$ を作り、fresh 変数 $\alpha_f$ に
\[
  C_s\subty\alpha_f
\]
を\emph{本体推論前に}登録する。その後
$\Gamma,f:(\alpha_f;\Xi_s)$ の下で $e_1$ を推論し、$C_{e_1}\subty\alpha_f$ を加える。
再帰中は単相、RHS 完了後に外側 level で一般化する。

\begin{intuition}
自己 skeleton の下界を先に置くことで、再帰呼び出しも外部呼び出しと同じ guard/paint 境界を通る。
本体を見た後で注釈を足す方式では、再帰呼び出しが境界を迂回してしまう。
\end{intuition}

\subsection{極性に基づく cleanup}

coalesced scheme $\tau$ 内で、id $a$ の guard/paint 出現の全体極性を調べる。
\[
\begin{aligned}
 \CovG_a(\tau)&\iff \text{$\gatom a{\Delta(a)}$ が共変位置に現れる},\\
 \CovP_a(\tau)&\iff \text{$\patom a$ が共変位置に現れる}.
\end{aligned}
\]
関数 domain で極性を反転し、computation result と row tail では保持する。

\begin{definition}[live id]
\[
  \Live(\tau)=\{a\mid \CovG_a(\tau)\land\CovP_a(\tau)\}.
\]
$\cleanup(\tau)$ は、まず各 occurrence の weight を正規化し、その後
$a\notin\Live(\tau)$ の guard/paint を全て消す。
\end{definition}

この定義は次の二段の議論を同時に表す。
\begin{enumerate}
  \item 共変 guard がなければ、共変 push を受け取ってそのまま外へ出す callback を注入できないため paint は観測不能である。
  \item 共変 paint がなければ、重みは scheme の外へ逃走できないため guard も観測不能である。
\end{enumerate}

\begin{example}[compose では消えない]
\[
 (\alpha\eff\emptyrow\to\beta\eff S)
 \to
 (\beta\eff\wocc S{\gatom a\Fam}\to\gamma\eff T)
 \to
 \alpha\eff\emptyrow\to\gamma\eff\wocc T{\patom a}
\]
では、$g$ の domain 内の guard は二回の反変を通って全体で共変であり、最終 result の paint も共変である。
したがって $a\in\Live(\tau)$ で、両者は残る。
\end{example}

\begin{conjecture}[cleanup の意味保存]
well-formed な推論結果 $\tau$ について、$a\notin\Live(\tau)$ の static markers と、対応して elaboration される
runtime color 操作を消去しても、任意の well-typed handler context における捕捉結果は変わらない。
\end{conjecture}

\section{型解決後の elaboration}

\subsection{なぜ後処理が必要か}

型推論中には、どの static id が cleanup 後に残るか確定しない。
そのため source term を直ちに runtime color 操作へ変換せず、まず静的 marker を持つ中間項 $M^\circ$ を生成する。
solver、coalescing、cleanup の後、残った weight に従って core term $M$ を作る。

\subsection{Core term}

必要な primitive を次とする。
\[
\begin{array}{rcl}
 M &::=& \cdots
 \mid \mathsf{newcolor}\ c\ \mathsf{in}\ M
 \mid \guard(c,H,V)
 \mid \pushid(c);M\\
 &&\mid\ \popid(c);M
 \mid \paint(c,V)
 \mid \scope(c,M).
\end{array}
\]
\[
  \mathsf{newcolor}\ c\ \mathsf{in}\ M
  \quad\equiv\quad
  \mathsf{let}\ c=\getid()\ \mathsf{in}\ M.
\]
$\scope(c,M)$ は説明上の derived form で、関数起動時に $c$ を push し、
通常 return でも effect escape でも、脱出時に $c$ を pop して外へ出る outcome 全体を paint する。

\begin{definition}[pop-paint]
\[
\begin{aligned}
 \poppaint(c,\ret V)
   &= \popid(c);\ret\paint(c,V),\\
 \poppaint(c,\oper(h,\colors,V,K))
   &= \popid(c);\oper(h,\colors,\paint(c,V),\paint(c,K)).
\end{aligned}
\]
operation payload と continuation も値なので paint の対象である。
\end{definition}

\subsection{Static id と dynamic id の一致}

elaboration environment
\[
  \theta:\Id\rightharpoonup\Color
\]
を使う。retained static id $a$ を instantiate した場所で
\[
  \mathsf{newcolor}\ c_a\ \mathsf{in}\ M
\]
を挿入し、$\theta(a)=c_a$ とする。同じ $a$ の
\[
 \guard(c_a,\Delta(a),-),\qquad
 \pushid(c_a),\qquad
 \poppaint(c_a,-)
\]
は必ずこの同じ binder を参照する。

\begin{formalbox}[title={Elaboration correspondence},fonttitle=\bfseries]
\begin{center}
\begin{tabular}{>{\raggedright\arraybackslash}p{0.25\linewidth}p{0.31\linewidth}p{0.31\linewidth}}
\toprule
静的情報 & 挿入する core marker & 動的意味 \\
\midrule
fresh static id $a$ & $c_a=\getid()$ & fresh color id \\
$\gatom aH$ & $\guard(c_a,H,V)$ & $H$ 外の visible effect を $c_a$ で着色 \\
関数起動 & $\pushid(c_a)$ & color flow に $c_a$ を登録 \\
$\patom a$ / 関数脱出 & $\poppaint(c_a,O)$ & flow から除き、outcome に paint \\
cleanup で $a$ 消去 & 対応 marker を消去 & 観測不能な color 機構を省略 \\
\bottomrule
\end{tabular}
\end{center}
\end{formalbox}

\subsection{guarded function の実行}

$\guard(c,H,V)$ が関数値なら、その適用は概念的に
\[
  \scope(c,\Guard_{c,H}(V\,W))
\]
へ展開される。したがって guard wrapper が保持する $c$ と、activation stack に push される $c$ は一致する。
$\Guard_{c,H}$ は次節で outcome transformer として定義する。

\begin{conjecture}[elaboration の一意性]
正規化・cleanup 済みの weighted scheme と、各 occurrence の provenance が与えられれば、
alpha-renaming と観測不能 marker の配置を除いて、$\getid$/guard/push/pop-paint の elaboration は一意である。
\end{conjecture}

\section{動的 color 意味論}

\subsection{Color flow と colored outcome}

runtime color を $c,d\in\Color$ とする。動的 color flow は stack
\[
  \flow ::= \epsilon \mid \flow\cdot c
\]
である。effect operation は color 列 $\colors$ を持つ。
\[
 O ::= \ret V \mid \oper(h,\colors,V,K).
\]
$K$ は continuation value である。

\begin{definition}[visibility]
\[
  \visible(\colors,\flow)
  \quad\Longleftrightarrow\quad
  \mathrm{set}(\colors)\cap\mathrm{set}(\flow)=\emptyset.
\]
handler は、その operation の color が現在の color flow に一つも含まれないときだけ捕捉する。
\end{definition}

\subsection{guard transformer}

ユーザの意図に合わせ、$\guard(c,H,V)$ は
\begin{quote}
color flow に含まれる color を既に持たない operation のうち、family が $H$ に含まれないものへ $c$ を塗る
\end{quote}
wrapper とする。outcome 上の transformer $\Guard_{c,H}^{\flow}$ を次で定める。
\[
\begin{aligned}
 \Guard_{c,H}^{\flow}(\ret V)
   &= \ret\guard(c,H,V),\\
 \Guard_{c,H}^{\flow}(\oper(h,\colors,V,K))
   &= \oper(h,\colors',\guard(c,H,V),
      \lambda x.\Guard_{c,H}^{\flow}(Kx)),
\end{aligned}
\]
ただし
\[
 \colors'=
 \begin{cases}
   c::\colors,
     & h\notin H\ \land\ \visible(\colors,\flow),\\
   \colors,&\text{otherwise}.
 \end{cases}
\]
$H$ 内の operation は無着色のまま見えるため、現在の handler が処理できる。
$H$ 外の visible operation は $c$ を得る。$c$ が flow に active なら、その operation は同じ flow 内の handler から見えなくなる。

\subsection{paint と scope}

paint は $H=\emptyset$ の guard と同じ「すべての visible effect へ color を残す」振る舞いを持つ。
\[
  \Paint_c^{\flow}=\Guard_{c,\emptyset}^{\flow}.
\]

$\scope(c,M)$ は $M$ を $\flow\cdot c$ の下で評価し、境界から出るときに $c$ を flow から外した後、outcome 全体へ $\Paint_c^{\flow}$ を適用する。
大ステップで概略を書くと、
\[
\inferrule*[right=E-Scope]
 {\flow\cdot c\vdash M\Downarrow O}
 {\flow\vdash \scope(c,M)\Downarrow \Paint_c^{\flow}(O)}.
\]
この一規則が、通常 return と effect による非局所脱出の双方で pop/paint を行う。

\subsection{Handler semantics}

$\handle_h M$ は $M$ の outcome を調べる。
\begin{itemize}
  \item $\ret V$ なら return branch を実行する。
  \item $\oper(h,\colors,V,K)$ かつ $\visible(\colors,\flow)$ なら operation branch を実行する。
  \item family が異なるか、$\neg\visible(\colors,\flow)$ なら operation を外へ再送し、continuation の後ろに同じ handler を付ける。
\end{itemize}
各 family に operation が一つなので、該当 family を捕捉した時点で incomplete branch はない。

\begin{lemma}[動的 guard hygiene]
$c\in\flow$ とする。$\guard(c,H,V)$ を通じて生じた family $h\notin H$ の visible operation は、
同じ $\flow$ の内側にある handler から捕捉されない。
\end{lemma}
\begin{proof}
operation が元々 flow color を持たないなら、guard transformer は color 列へ $c$ を追加する。
$c\in\flow$ なので新しい color 列と flow の共通部分は空でなく、visibility 条件が偽になる。
元々 flow color を持つなら、最初から visibility が偽である。いずれの場合も handler の捕捉規則を適用できない。
\end{proof}

\begin{corollary}[許可 family の可視性]
$h\in H$ で、operation の既存 color が flow と交わらないなら、$\guard(c,H,-)$ はその color を変更しないため、
family $h$ の handler は捕捉できる。
\end{corollary}

\begin{lemma}[fresh id separation]
異なる scheme instantiation が $\getid$ により fresh な $c_1\neq c_2$ を得るなら、片方の push/pop は他方の guard visibility を直接変えない。
\end{lemma}
\begin{proof}
visibility は color の一致で判定される。$c_1\neq c_2$ なので、一方だけを flow に追加・削除しても、他方との集合交差の真偽は変わらない。
\end{proof}

\section{静的重みと動的 color の対応}

静的 guard は guard wrapper と activation push の組、静的 paint は scope exit の pop-paint を抽象化する。
一つの値の流れに沿って wrapper を合成すると、同じ id の guard/paint 回数の差だけが型上の係数として残る。

\begin{definition}[weight abstraction]
動的 provenance path $\pi$ に含まれる id $a$ の guard wrapper 数を $G_a(\pi)$、pop-paint 数を $P_a(\pi)$ とし、
\[
  \mathsf{abs}(\pi)(a)=G_a(\pi)-P_a(\pi)
\]
とする。
\end{definition}

\begin{lemma}[path composition]
\[
  \mathsf{abs}(\pi_1\cdot\pi_2)=\mathsf{abs}(\pi_1)+\mathsf{abs}(\pi_2).
\]
\end{lemma}
\begin{proof}
連結 path 上の guard 数と paint 数はいずれも各部分 path の和である。差を取れば点ごとの重み加算になる。
\end{proof}

この補題が、weighted inequality の replay で重みを加算する動的直感を与える。
ただし動的 path の順序正当性は runtime stack が保証し、静的加算だけから任意の push/pop 順序を許すわけではない。

\begin{conjecture}[静的・動的対応]
型推論と elaboration が成功した source term $e$ について、elaborated core term の任意の値 provenance path $\pi$ に現れる
color 操作の abstraction は、その値出現の weighted type に付く正規化済み weight と一致する。
\end{conjecture}

\begin{conjecture}[型保存と handler hygiene]
$\Gamma\vdash e\Rightarrow C$ が成功し、$e$ を型指向 elaboration して $M$ を得たとする。
$M$ の評価で、guard capability が許していない family を同じ dynamic scope 内の handler が捕捉することはない。
また評価が一歩進んでも computation type $C$ と weighted bounds の充足可能性は保存される。
\end{conjecture}

\section{主型性への証明計画}

\subsection{制約生成の主性}

ordinary Simple-sub の証明を次の点で拡張する。
\begin{enumerate}
  \item weighted edge は ordinary subtype edge に有限台整数ラベルを付けただけで、構造分解は決定的である。
  \item bounds replay のラベルは加法で一意に合成される。
  \item weighted row split は fresh tail を一つ導入し、head と tail を直接混ぜないため、最一般 residual を保つ。
  \item annotation は ordinary row を上から閉じず、weight exposure だけを検査するので、自由な row tail を失わない。
  \item let generalization は型変数・row 変数・id を同時に freshen する。
\end{enumerate}

\begin{conjecture}[Principal weighted scheme]
有限 family signature、完全 handler、disjoint row head/tail、monomorphic recursion の下で、
成功する項 $e$ の constraint generation と weighted Simple-sub solver は principal weighted scheme を生成する。
任意の別の well-typed scheme は、ordinary type/row 変数の substitution、static id の injective renaming、
および許される subtyping により principal scheme から得られる。
\end{conjecture}

\subsection{停止性}

型対と正規化済み左右 weight の cache を使う。id の数は構文と instantiate 回数に対して有限であり、
各比較 path の weight は finite-support integer vector である。ただし再帰 bounds replay で係数が無限に増える可能性を除くため、
well-bracketed elaboration 由来の balance invariant、または SCC ごとの飽和・hash-cons 条件が必要である。

\begin{conjecture}[Solver termination]
well-bracketed annotation translation が生成する制約に限定し、同一 SCC 内の同一 id について guard/paint balance が有界であるなら、
weighted constrain は有限回で停止する。
\end{conjecture}

\begin{statusbox}
停止性はこの稿で最も未確定な箇所である。signed coefficient を無制限に許すと、再帰 replay が同じ id の係数を増やし続ける可能性がある。
実装では飽和、SCC 正規化、または「同じ activation id を再帰周回で追加発行しない」という生成不変条件のいずれかを明文化する必要がある。
\end{statusbox}

\section{設計上の決定と残る論点}

\subsection{この稿で固定したこと}
\begin{itemize}
  \item 値型は推論し、ユーザ注釈は effect slot だけを書く。
  \item 型変数そのものではなく参照出現に重みが付く。
  \item 重みは左右付き不等式と bounds replay を通る。
  \item handler split は必ず fresh tail を挟む。
  \item 各 family は一 operation、handler は完全。
  \item let scheme は static id も量化し、instantiate ごとに freshen する。
  \item 再帰は自己 skeleton 下界を本体より前に置く。
  \item 静的 guard/paint と動的 push/pop を別物として定義する。
  \item dynamic guard は $H$ 外の visible effect を着色し、handler は flow にない color の effect を捕捉する。
  \item 通常 return と effect escape の双方で pop/paint を行う。
\end{itemize}

\subsection{次に議論すべき点}
\begin{enumerate}
  \item recursive replay における signed coefficient の有界性を、生成規則だけで証明できるか。
  \item $\guard(c,H,V)$ が record、tuple、closure、continuation など複合値へどう分配されるか。
  \item effect family に型引数を戻したとき、family intersection と不変引数制約をどこで生成するか。
  \item cleanup の $\Live$ 条件が十分条件に留まるか、必要十分条件まで強められるか。
  \item elaboration marker の配置一意性と、cleanup 後 marker 消去の文脈同値性。
  \item deep handler、multi-shot continuation、並行評価を入れたときの color flow の複製規則。
\end{enumerate}

\section{まとめ}

本稿の体系では、effect-only 注釈が static id を発行し、値の\emph{参照出現}に guard weight を付ける。
型比較は左右の重みを保ったまま構造分解され、Simple-sub の lower/upper bounds replay で重みが加算される。
handler は row を直接 weighted tail へ書き換えず、ordinary upper bound と fresh tail を作った後、tail 間に weighted edge を張る。
let 多相では id も量化し、instantiate ごとに fresh 化する。一般化時には極性を見て、共変 guard と共変 paint の双方を持つ id だけを残す。

型解決後の elaboration は static id を dynamic color id へ具体化する。guard と push は同じ $\getid$ の結果を共有し、
関数脱出では pop と paint を一緒に実行する。paint は return value だけでなく、operation payload と continuation にも付着する。
動的 handler は現在の color flow に含まれない color の operation だけを捕捉する。
このため、guard capability 外の effect は同じ scope 内の handler から隠され、許可された effect だけが見える。

静的には signed monoid、動的には本物の stack と粘着する color、という二層分離により、
Simple-sub の主型推論と runtime hygiene の双方を一つの説明へ接続できる。

\begin{thebibliography}{9}
\bibitem{parreaux2020}
Lionel Parreaux.
\newblock The Simple Essence of Algebraic Subtyping: Principal Type Inference with Subtyping Made Easy.
\newblock \emph{Proceedings of the ACM on Programming Languages}, 4(ICFP), 2020.

\bibitem{workingnote}
Yulang effect subtraction working note.
\newblock \emph{effect subtraction（stack/filter 重み）}, 2026年5月31日版およびその後の議論。
\end{thebibliography}

\end{document}
